% BHU PYQ -- Manually transcribed & error-corrected
% Paper: BCF-211 : Corporate Accounting
% Exam: B.Com. (Hons.) F.M.M. Semester III Examination, 2023-24

\documentclass[11pt,a4paper]{article}
\usepackage[a4paper,top=2.5cm,bottom=2.5cm,left=2.8cm,right=2.8cm,headheight=14pt]{geometry}
\usepackage{amsmath,amssymb}
\usepackage[T1]{fontenc}
\usepackage[utf8]{inputenc}
\usepackage{lmodern,microtype,fancyhdr,enumitem,hyperref,lastpage,parskip,booktabs}

\hypersetup{
    colorlinks=true,
    urlcolor=black,
    linkcolor=black,
    pdftitle={BCF-211: Corporate Accounting -- BHU B.Com. (Hons.) FMM Sem III 2023-24}
}

\pagestyle{fancy}
\fancyhf{}
\renewcommand{\headrulewidth}{0.4pt}
\renewcommand{\footrulewidth}{0.4pt}
\fancyhead[L]{\small\textit{Banaras Hindu University}}
\fancyhead[R]{\small\textit{BCF-211: Corporate Accounting}}
\fancyfoot[C]{\small Page \thepage\ of \pageref{LastPage}}

\newlist{parts}{enumerate}{2}
\setlist[parts,1]{label=(\alph*),leftmargin=*,itemsep=5pt,topsep=3pt}
\setlist[parts,2]{label=(\roman*),leftmargin=*,itemsep=3pt,topsep=2pt}
\newcommand{\pts}[1]{\hfill\textit{[#1]}}

\begin{document}

\begin{center}
  {\large\bfseries BANARAS HINDU UNIVERSITY}\\[2pt]
  {\normalsize B.Com. (Hons.) F.M.M. --- Semester III Examination, 2023-24}\\[6pt]
  \rule{0.8\linewidth}{0.5pt}\\[6pt]
  {\large\bfseries Paper No. BCF-211: Corporate Accounting}\\[4pt]
  {\normalsize Time: 3 Hours\hfill Maximum Marks: 70}
\end{center}

\rule{\linewidth}{0.4pt}

\smallskip
\noindent\textit{Instructions: Answer all questions. All questions carry equal marks (14 marks each).}

\bigskip

\noindent\textbf{Question 1.}
`S' India Ltd.\ invited applications for 20,000 equity shares of Rs.\ 100 each, payable Rs.\ 40 on application, Rs.\ 30 on allotment and Rs.\ 30 on first and final call. The company received applications for 32,000 shares. Applications for 2,000 shares were rejected and money returned to applicants. Applications for 10,000 shares were accepted in full and applicants for 20,000 shares allotted half of the number of shares applied and excess application money adjusted into allotment. All money due on allotment and call was received.

Pass journal entries and prepare Cash Book.\pts{14}

\centerline{\textbf{OR}}

What is underwriting of shares? What are the main problems of underwriting of shares? Who can be an underwriter of shares?\pts{14}

\medskip\rule{\linewidth}{0.2pt}

\noindent\textbf{Question 2.}
What is the theory of amalgamation of companies? Discuss its main objectives and the methods of amalgamation.\pts{14}

\centerline{\textbf{OR}}

The following are the Balance Sheets of two firms M/s.\ A \& B and M/s.\ X \& Y as on 31st March, 2023:\pts{14}

\begin{center}\small
  \begin{tabular}{p{2.6cm}rr|p{2.6cm}rr}
    \toprule
    \textbf{Liabilities} & \textbf{A \& B} & \textbf{X \& Y} & \textbf{Assets} & \textbf{A \& B} & \textbf{X \& Y} \\
    & \textbf{(Rs.)} & \textbf{(Rs.)} & & \textbf{(Rs.)} & \textbf{(Rs.)} \\
    \midrule
    Sundry Creditors & 30,000 & 25,000 & Cash at Bank & 5,600 & 6,700 \\
    Mrs.\ A's Loan & 20,000 & -- & Stock & 20,400 & 18,300 \\
    Capital A/cs: & & & Debtors & 10,000 & 20,000 \\
    \quad A & 15,000 & -- & Investment & -- & 15,000 \\
    \quad B & 15,000 & -- & Furniture & 4,000 & 5,000 \\
    \quad X & -- & 24,000 & Premises & 40,000 & -- \\
    \quad Y & -- & 16,000 & & & \\
    \midrule
    \textbf{Total} & \textbf{80,000} & \textbf{65,000} & \textbf{Total} & \textbf{80,000} & \textbf{65,000} \\
    \bottomrule
  \end{tabular}
\end{center}

The two firms decided to amalgamate their businesses as from 1st April, 2023. For this purpose, it was agreed that Mrs.\ A's Loan should be repaid and that the investments of M/s.\ X and Y be not taken over by the new firm. Premises were revalued at Rs.\ 50,000 but the stock of M/s.\ A and B was found overvalued by Rs.\ 4,000.

The stock of M/s.\ X and Y was undervalued by Rs.\ 2,000. A provision of 5\% was created for bad debts of both the firms. The total capital of the new firm was to be Rs.\ 80,000 and the capital of each partner was to be in his profit-sharing ratio which was to be 3:2:3:2 respectively.

Close the books of the two firms and pass opening entries in the books of M/s.\ A, B, and Y. Also prepare the balance sheet of the newly constituted firm.

\medskip\rule{\linewidth}{0.2pt}

\noindent\textbf{Question 3.}
What is meant by reduction of capital? Give the journal entries passed in the books of accounts for reduction of capital.\pts{14}

\centerline{\textbf{OR}}

The share capital of Z Ltd.\ consisted of the following:\pts{14}
\begin{parts}
  \item 10,000 6\% preference shares of Rs.\ 10 each and
  \item 50,000 equity shares of Rs.\ 10 each
\end{parts}
The shares were fully paid. By the end of the year, it had accumulated losses of Rs.\ 3,50,000 besides preliminary expenses of Rs.\ 20,000. Fixed assets which stood in the books at Rs.\ 14,00,000 were overvalued to the extent of Rs.\ 4,00,000. A scheme of capital reduction was adopted and approved. Under the scheme, 6\% preference shares were to be converted into 7\% preference shares of Rs.\ 6 each and equity shares to be converted into shares of Rs.\ 2 each. The preference shares dividend which were in arrears for three years were to be cancelled.

Pass necessary journal entries in the books of the company.

\medskip\rule{\linewidth}{0.2pt}

\noindent\textbf{Question 4.}
Explain the Liquidator's Final Statement of Accounts. Give the format with imaginary figures.\pts{14}

\centerline{\textbf{OR}}

`A' Ltd.\ went into liquidation on 31-03-2022 with the following liabilities:\pts{14}
\begin{parts}
  \item Secured Creditors Rs.\ 2,00,000 (securities realized Rs.\ 2,50,000)
  \item Preferential Creditors Rs.\ 6,000
  \item Unsecured Creditors Rs.\ 3,05,000
\end{parts}
The liquidator met liquidation expenses amounting to Rs.\ 2,520. The liquidator is entitled to remuneration @ 3\% on amount realized including secured assets held by secured creditors and 1.5\% on amount distributed to unsecured creditors. Assets (other than secured assets realized) realized Rs.\ 2,60,000.

Prepare Liquidator's Final Statement of Account.

\medskip\rule{\linewidth}{0.2pt}

\noindent\textbf{Question 5.}
Define holding and subsidiary company. What is a consolidated balance sheet and how is it prepared?\pts{14}

\centerline{\textbf{OR}}

`H' Ltd.\ acquired all the shares of `S' Ltd.\ on 1st Jan., 2022. The Balance Sheet of the two companies on 31st March, 2022 is as under:\pts{14}

\begin{center}\small
  \begin{tabular}{p{5.5cm}rr}
    \toprule
    \textbf{Particulars} & \textbf{H Ltd. (Rs.)} & \textbf{S Ltd. (Rs.)} \\
    \midrule
    \textbf{I. Equity \& Liabilities} & & \\
    Shareholders' Funds & & \\
    \quad Share Capital & 50,000 & 30,000 \\
    Reserves and Surplus & & \\
    \quad General Reserve (01/04/2021) & 20,000 & 15,000 \\
    \quad Statement of Profit \& Loss (Surplus) & 25,000 & 10,000 \\
    Current Liabilities & & \\
    \quad Trade Payables & 20,000 & 15,000 \\
    \midrule
    \textbf{Total} & \textbf{1,15,000} & \textbf{70,000} \\
    \midrule
    \textbf{II. Assets} & & \\
    Sundry Assets & 65,000 & 70,000 \\
    Shares in S Ltd. at cost & 50,000 & -- \\
    \midrule
    \textbf{Total} & \textbf{1,15,000} & \textbf{70,000} \\
    \bottomrule
  \end{tabular}
\end{center}

The Profit \& Loss of S Ltd.\ had a credit balance of Rs.\ 3,000 on 01/04/2021. The Profit of S Ltd.\ accrued evenly throughout the year. Prepare a consolidated Balance Sheet.

\vfill
\rule{\linewidth}{0.4pt}
\begin{center}\small
\href{https://bhu-exam.vercel.app/}{Click here to visit our website}\\[4pt]
\href{https://me-aryan.vercel.app/}{Click here to visit our contributor}\\[4pt]
\href{https://quashan-me.vercel.app/}{Click here to visit our contributor}
\end{center}

\end{document}
